% LSPV-5 REPORT::DUALPRIMERET1::2026-10-05T07:08:00-07:00
\documentclass[11pt]{article}
\usepackage[margin=.86in]{geometry}
\usepackage{amsmath,amssymb}
\begin{document}
\section*{DUALPRIMERET1: dual-prime literal fingerprints as a SAT-semantic return candidate}
\textbf{STATUS: PROVED/EXACT local encoding lemmas + COMPUTATIONAL controls + CANDIDATE return theorem. P versus NP remains OPEN.}

\section*{Construction}
For each Boolean variable x_i choose distinct primes t_i,f_i. For assignment a define
\[
Z(a)=\prod_{i:a_i=1}t_i\prod_{i:a_i=0}f_i.
\]
Then a is recovered exactly by divisibility: t_i divides Z(a) iff a_i=1 and f_i divides Z(a) iff a_i=0. The user's example x_1 true =3 and x_1 false =5 is a valid instance of this scheme.

For a literal ell define p(ell)=t_i for ell=x_i and p(ell)=f_i for ell=not x_i. A clause C is satisfied by a exactly when
\[
\gcd\left(Z(a),\prod_{\ell\in C}p(\ell)\right)>1.
\]
Therefore a CNF F is satisfiable exactly when there exists a legal transversal product Z choosing one prime from each pair such that this gcd condition holds for every clause.

\section*{Exact theorem: DUALPRIME-SAT}
Let P_i={t_i,f_i} be pairwise disjoint prime pairs and B_C=product over satisfying-literal primes of clause C. Then
\[
F\in SAT
\Longleftrightarrow
\exists Z\;[
\forall i,\ |\{t_i,f_i\}\cap\operatorname{PF}(Z)|=1
\ \wedge\
\forall C,\gcd(Z,B_C)>1].
\]
Proof is the direct assignment-to-transversal bijection. This is an exact SAT-semantic arithmetic representation, not yet a smaller ordinary-CNF return.

\section*{Information persistence}
Multiplication is useful because prime support is union under multiplication and exact factor support can be recovered from a square-free product. Thus a sequence of operations that only multiplies selected literal/event primes can preserve a lossless commutative fingerprint. However, order and alternative derivations are not preserved by a plain product, and Boolean negation is not represented as multiplication. PRIMECRT1 already proves these limitations and supplies CRT idempotents when exact Boolean composition is needed.

The important correction is that later factorization must not be treated as the hard operation. The generating computation knows the factors, and integer gcd/divisibility already reads the literal conditions without full factorization. The possible value is therefore \emph{semantic persistence}, not factoring hardness.

\section*{Candidate SAT-return theorem}
A potentially useful new target is the following.

\textbf{DUALPRIME-RETURN, CANDIDATE.}
Suppose a source transformation maps F to a polynomial-length collection of arithmetic constraints
\[
\mathcal D(F)=\{(P_i),(B_C),\mathcal O\}
\]
where \mathcal O consists only of source-certified operations preserving prime-support semantics. If one can prove a uniform compiler Return such that
\[
F\in SAT\iff Return(\mathcal D(F))\in SAT,
\]
the returned object is an explicit ordinary CNF, its length is T^{zeta+o(1)} with zeta<1 or belongs to a recursively contractive typed multi-return family, and all construction/acquisition costs are subcritical, then DUALPRIME-RETURN instantiates the missing SAT-semantic return edge.

No such compiler is proved here.

\section*{Hostile audit}
\textbf{Semantic quotient.} The prime transversal representation is polynomially interconvertible with the ordinary assignment representation. It does not create new SAT information.

\textbf{TRACECLOSE1 / PRIME-FNP-CLOSE1.} Under P=NP, any polynomially balanced NP-verifiable legal prime-labelled witness can be selected in FP. Deterministic polynomial trace labels can be regenerated. Hence provenance alone cannot be the contradiction.

\textbf{Factoring.} Factoring is not chargeable as hardness. GCD/divisibility suffices for clause tests, and the encoder knows the prime decomposition.

\textbf{Commutativity.} Products forget operation order and multiplicity unless positional/event-specific primes are added. Such labels remain compact and replayable in polynomial traces.

\textbf{Algebrization.} The construction is explicitly arithmetic and survives algebraic-oracle replacement; it therefore does not itself supply the nonalgebrizing hinge.

\textbf{Easy controls.} Horn, 2-SAT, affine/Tseitin, bounded-treewidth and independent-product formulas admit the same encoding. Large prime fingerprints therefore cannot be read as hardness.

\textbf{Return typing.} The arithmetic object is not an ordinary-CNF query. A model-matched compiler to a smaller CNF remains the load-bearing OPEN arrow.

\section*{COMPUTATIONAL test}
Committed experiment experiments/return_map/dualprime1.py exhaustively checks n=5 assignment roundtrip and exact clause/formula divisibility semantics for a sample 3-CNF. This is a finite semantics test only.

\section*{Relation to SIL/BOOL}
In SIL terms, the dual-prime product is a realization-level annotation of a boundary/witness relation. Projection or composition may preserve a prime-support summary only when an explicit transport rule is supplied. In BOOL terms, each transformation must state its preservation signature: SAT, models/witness recoverability, prime-support relation, representation size, and construction cost. SAME_REALIZATION requires a CompilationMap rather than equality of numerical labels after arbitrary preprocessing.

This suggests a stronger experiment than raw provenance: search for a source transformation whose \emph{prime-support invariant is sufficient for an exact SAT combiner and contracts under projection}. If the invariant is merely reconstructible metadata, TRACECLOSE1/PRIME-FNP-CLOSE1 closes it. If it yields an explicit smaller SAT query without current-scale reacquisition, it becomes a genuine return candidate.

\section*{Next test: PRIMEPROJ1}
For small 3-CNFs, carry dual-prime support through SIL restrict, project, join and compose. At each step compute the minimal exact boundary relation and compare:
(1) full dual-prime support,
(2) the smallest prime-support quotient sufficient for every future SAT context,
(3) explicit CNF size of that quotient,
(4) whether the quotient can be obtained without solving the original SAT instance.
The desired phenomenon is a recursively contractive SAT-sufficient quotient. Failure by exponential contextual index, full-input carry, or SAT-dependent acquisition should be recorded immediately.

\section*{Status}
DUALPRIME-SAT is PROVED/EXACT as an arithmetic encoding identity. Prime-support persistence under multiplication is PROVED/EXACT. Prime factorization as a hardness resource remains FALSIFIED/RETIRED. DUALPRIME-RETURN and a contractive SIL prime-support quotient are CANDIDATE/OPEN. Arbitrary-algorithm acquisition, ordinary-CNF return, nonrelativization, nonalgebrization, and P versus NP remain OPEN.

\end{document}
% END_OF_VERIFIED_REPORT::LSPV-5::DUALPRIMERET1
